\SourceLeaf{13}
\begin{center}\Large A NEW PROOF\end{center}

\begin{center}\large OF AN ARITHMETICAL THEOREM.\end{center}

\BookRule{35mm}

\Article{1}

Questions in higher arithmetic very often present a singular phenomenon, which occurs much more rarely in analysis and contributes greatly to their charm. In analytic investigations we usually cannot reach new truths unless we have first thoroughly mastered the principles on which they rest and which must, as it were, open the way to them. In arithmetic, by contrast, most elegant new truths very frequently spring forth through induction and a kind of unexpected good fortune; their proofs lie so deeply hidden and are wrapped in such darkness that they elude every attempt and deny access to the keenest investigations. Moreover, so great and so wonderful is the connection among arithmetical truths that at first sight seem utterly different, that quite often, while we are seeking something quite different, we at last most happily reach the proof so greatly desired and previously sought in vain through long reflection, by a path far different from the one expected. Usually, however, truths of this kind are of such a nature that they can be approached by several very different paths, and the paths that first present themselves are not always the shortest. It must therefore certainly be highly valued if, after such a truth has long been investigated in vain, and has then indeed been proved, but by rather obscure detours, we finally succeed in discovering the simplest and genuine path.

\Article{2}

Among the questions discussed in the preceding article, a distinguished place is held by the theorem containing almost the whole theory of quadratic residues, which in the \emph{Arithmetical Investigations} (Section IV) is distinguished by the name of the \emph{fundamental theorem}\SourceLeaf{14}. The illustrious \textsc{Legendre} must undoubtedly be regarded as the \emph{first} discoverer of this most elegant theorem, after the great geometers \textsc{Euler} and \textsc{Lagrange} had, long before, already discovered several of its special cases by induction. I do not pause here to enumerate these men's attempts at a proof; those who wish may consult the work just mentioned. Let me add only, in confirmation of what was said in the preceding article, what concerns my own attempts. I came upon the theorem itself independently in the year 1795, when I was entirely ignorant of everything that had already been worked out in higher arithmetic and wholly cut off from access to the literature. But for an entire year it tormented me and eluded my most strenuous efforts, until at last I obtained the proof given in the fourth section of that work. Afterwards three others, resting on wholly different principles, presented themselves to me; I gave one in the fifth section, and the others, not inferior to it in elegance, I shall make public on another occasion. But all these proofs, although they seem to leave nothing to be desired in rigour, are derived from principles that are too different in kind, perhaps excepting the first, which nevertheless proceeds by more laborious reasoning and is burdened by rather lengthy operations. I therefore do not hesitate to declare that a \emph{genuine} proof has so far been lacking. Let the experts now judge whether the one that I recently succeeded in discovering, and that the following pages present, deserves to be honoured with this name.

\Article{3}

\textsc{Theorem.} \emph{Let $p$ be a positive prime number; $k$ any integer not divisible by $p$;} \emph{use the following definitions:}

\[
\begin{aligned}
A &\quad \text{\emph{collection of the numbers}}\quad 1,\ 2,\ 3\ \dotsfour\ \tfrac12(p-1)\\
B &\quad \text{\emph{collection of these}}\quad \tfrac12(p+1),\ \tfrac12(p+3),\ \tfrac12(p+5)\ \dotsfour\ p-1
\end{aligned}
\]

\emph{Take the products of $k$ by the individual numbers in $A$, and their least positive residues modulo $p$, which are evidently all distinct and belong partly to $A$, partly to} \emph{$B$. Now if altogether $B$ contains $\mu$ of the residues, $k$ will be either a quadratic residue or a quadratic nonresidue of $p$, according as $\mu$ is even or odd.}

\emph{Proof.} Let the residues belonging to $A$ be $a,\ a'\ a''\ \dotsfour$, and the remaining ones, belonging to $B$, be $b,b',b''\dotsfour$. It is clear that the complements of the latter, $p-b,p-b',p-b''\dotsfour$, are all distinct from the numbers $a,\ a',\ a''\ \dotsfour$, but together with these fill out the collection\SourceLeaf{15} $A$. Thus we have

\[
\begin{aligned}
1\pdot2\pdot3\ \dotsfour\ \tfrac12(p-1)=aa' a''\ \dotsfour\ (p-b)(p-b')(p-b'')\ \dotsfour
\end{aligned}
\]

But the latter product evidently becomes

\[
\begin{aligned}
&\equiv(-1)^\mu aa' a''\ \dotsfour\ bb' b''\ \dotsfour\ \equiv(-1)^\mu k\pdot2k\pdot3k\ \dotsfour\ \tfrac12(p-1)k\\
&\equiv(-1)^\mu k^{\frac12(p-1)}\,1\pdot2\pdot3\ \dotsfour\ \tfrac12(p-1) (\mathrm{mod.}p)
\end{aligned}
\]

Hence we have

\[
\begin{aligned}
1\equiv(-1)^\mu k^{\frac12(p-1)}
\end{aligned}
\]

or $k^{\frac12(p-1)}\equiv\pm1$, according as $\mu$ is even or odd; our theorem follows immediately.

\Article{4}

We can greatly shorten the following reasoning by introducing some suitable notation. Thus the symbol $(k,p)$ will denote for us the number of products among

\[
\begin{aligned}
k,\ 2k,\ 3k\ \dotsfour\ \tfrac12(p-1)k,
\end{aligned}
\]

whose least positive residues modulo $p$ exceed half of it. Further, when $x$ is any noninteger quantity, we shall use the symbol $[x]$ to denote the integer immediately smaller than $x$ itself, so that $x-[x]$ is always a positive quantity lying between the limits $0$ and $1$. The following relations are now obtained with little trouble:

I. $[x]+[-x]=-1.$

II. $[x]+h=[x+h]$, whenever $h$ is an integer.

III. $[x]+[h-x]=h-1.$

IV. If $x-[x]$ is a fraction smaller than $\tfrac12$, then $[2x]-2[x]=0$; but if $x-[x]$ is greater than $\tfrac12$, then $[2x]-2[x]=1.$

V. Thus, when the least positive residue of the integer $h$ modulo $p$ lies below $\tfrac12p$, we have $\gbr{\frac{2h}{p}}-2\gbr{\frac{h}{p}}=0$; but when that residue lies above $\tfrac12p$, we have $\gbr{\frac{2h}{p}}-2\gbr{\frac{h}{p}}=1.$

\SourceLeaf{16}
VI. It follows at once that $(k,p)=$

\[
\begin{aligned}
&\gbr{\frac{2k}{p}}+\gbr{\frac{4k}{p}}+\gbr{\frac{6k}{p}}\dotsfour+\gbr{\frac{(p-1)k}{p}}\\
&-2\gbr{\frac{k}{p}}-2\gbr{\frac{2k}{p}}-2\gbr{\frac{3k}{p}}\dotsfour-2\gbr{\frac{\frac12(p-1)k}{p}}.
\end{aligned}
\]

VII. From VI and I we obtain without difficulty

\[
\begin{aligned}
(k,p)+(-k,p)=\tfrac12(p-1)
\end{aligned}
\]

It follows that $-k$ has either the same or the opposite relation to $p$ (insofar as it is a quadratic residue or nonresidue of it) as $+k$, according as $p$ is of the form $4n+1$ or of the form $4n+3$. In the former case $-1$ is evidently a residue, in the latter a nonresidue, of $p$.

\Needspace{10\baselineskip}VIII. We shall transform the formula given in VI as follows. By III we have

\[
\begin{aligned}
\gbr{\frac{(p-1)k}{p}}&=k-1-\gbr{\frac{k}{p}},\\
\gbr{\frac{(p-3)k}{p}}&=k-1-\gbr{\frac{3k}{p}},\\
\gbr{\frac{(p-5)k}{p}}&=k-1-\gbr{\frac{5k}{p}}\dotsfour
\end{aligned}
\]

Applying these substitutions to the last $\frac{p\mp1}{4}$ terms of the upper series in that expression, we have

\Needspace{8\baselineskip}\emph{first}, whenever $p$ is of the form $4n+1$

\[
\begin{aligned}
(k,p)&=\tfrac14(k-1)(p-1)\\
& -2\left\{\gbr{\frac{k}{p}}+\gbr{\frac{3k}{p}}+\gbr{\frac{5k}{p}}\dotsfour+\gbr{\frac{\frac12(p-3)k}{p}}\right\}\\
&\qquad-\left\{\gbr{\frac{k}{p}}+\gbr{\frac{2k}{p}}+\gbr{\frac{3k}{p}}\dotsfour+\gbr{\frac{\frac12(p-1)k}{p}}\right\}
\end{aligned}
\]

\Needspace{8\baselineskip}\emph{second}, whenever $p$ is of the form $4n+3$

\[
\begin{aligned}
(k,p)&=\tfrac14(k-1)(p+1)\\
& -2\left\{\gbr{\frac{k}{p}}+\gbr{\frac{3k}{p}}+\gbr{\frac{5k}{p}}\dotsfour+\gbr{\frac{\frac12(p-1)k}{p}}\right\}\\
&\qquad-\left\{\gbr{\frac{k}{p}}+\gbr{\frac{2k}{p}}+\gbr{\frac{3k}{p}}\dotsfour+\gbr{\frac{\frac12(p-1)k}{p}}\right\}
\end{aligned}
\]

IX. For the special case $k=+2$, the formulas just given yield $(2,p)=\tfrac14(p\mp1)$, taking the upper or lower sign according as $p$ is of the form $4n+1$ or $4n+3$. Thus $(2,p)$ is even, and hence $2Rp$, whenever $p$ is of the form $8n+1$ or $8n+7$; on the other hand, $(2,p)$ is odd and $2Np$, whenever $p$ is of the form $8n+3$ or $8n+5$.

\SourceLeaf{17}
\Needspace{14\baselineskip}\Article{5}

\textsc{Theorem.} \emph{Let $x$ be a positive noninteger quantity, among whose multiples $x$, $2x$,} \emph{$3x\dotsfour$ up to $nx$ none is an integer; put $[nx]=h$, from which it is easily concluded} \emph{that among the multiples of the reciprocal quantity $\frac1x$, $\frac2x$, $\frac3x\dotsfive$ up to $\frac hx$ an integer is likewise not} \emph{found. Then I say that we have}

\[
\begin{aligned}
\left.\begin{aligned}
&[x]+[2x]+[3x]\dotstwo+[nx]\\
&+\gbr{\frac1x}+\gbr{\frac2x}+\gbr{\frac3x}\dotsfour+\gbr{\frac hx}
\end{aligned}\right\}=nh.
\end{aligned}
\]

\emph{Proof.} In the series $[x]+[2x]+[3x]\dotsfour+[nx]$, which we shall put $=\Omega$, the first terms up to and including the term in position $\gbr{\frac1x}$ are evidently all $=0$; the following ones up to the term in position $\gbr{\frac2x}$ are all $=1$; the following ones up to the term in position $\gbr{\frac3x}$ are all $=2$, and so on. Hence we have

\[
\begin{aligned}
\left.\begin{aligned}
\Omega={}&0\times\gbr{\frac1x}\\
&+1\times\left\{\gbr{\frac2x}-\gbr{\frac1x}\right\}\\
&+2\times\left\{\gbr{\frac3x}-\gbr{\frac2x}\right\}\\
&+3\times\left\{\gbr{\frac4x}-\gbr{\frac3x}\right\}\\
&\hspace{12mm}\text{etc.}\\
&+(h-1)\left\{\gbr{\frac hx}-\gbr{\frac{h-1}{x}}\right\}\\
&+h\left\{n-\gbr{\frac hx}\right\}
\end{aligned}\right\}=hn-\gbr{\frac1x}-\gbr{\frac2x}-\gbr{\frac3x}\dotsfour-\gbr{\frac hx}
\end{aligned}
\]

Q. E. D.

\Needspace{12\baselineskip}\Article{6}

\textsc{Theorem.} \emph{If $k,p$ denote any positive odd numbers relatively prime to each other, we have}

\[
\begin{aligned}
\left.\begin{aligned}
&\gbr{\frac kp}+\gbr{\frac{2k}{p}}+\gbr{\frac{3k}{p}}\dotsfour+\gbr{\frac{\frac12(p-1)k}{p}}\\
&+\gbr{\frac pk}+\gbr{\frac{2p}{k}}+\gbr{\frac{3p}{k}}\dotsfour+\gbr{\frac{\frac12(k-1)p}{k}}
\end{aligned}\right\}=\tfrac14(k-1)(p-1).
\end{aligned}
\]

\emph{Proof.} Supposing, as we may, that $k<p$, we have $\frac{\frac12(p-1)k}{p}$ less than $\frac12 k$, but greater than $\frac12(k-1)$, and therefore $\gbr{\frac{\frac12(p-1)k}{p}}=\frac12(k-1)$. Thus the present theorem evidently follows immediately from the preceding one by putting there $\frac kp=x$, $\frac12(p-1)=n$, and hence $\frac12(k-1)=h$.

\SourceLeaf{18}
\Needspace{8\baselineskip}Moreover, it can be proved in a similar manner that, if $k$ is an \emph{even} number relatively prime to $p$, we have

\[
\begin{aligned}
\left.\begin{aligned}
&\gbr{\frac kp}+\gbr{\frac{2k}{p}}+\gbr{\frac{3k}{p}}\dotsfour+\gbr{\frac{\frac12(p-1)k}{p}}\\
&+\gbr{\frac pk}+\gbr{\frac{2p}{k}}+\gbr{\frac{3p}{k}}\dotsfour+\gbr{\frac{\frac12 kp}{k}}
\end{aligned}\right\}=\tfrac14 k(p-1)
\end{aligned}
\]

But we do not dwell on this proposition, which is not needed for our purpose.

\clearpage\Article{7}

The fundamental theorem now follows immediately by combining the preceding theorem with proposition VIII of Article 4. Namely, if $k,p$ denote any unequal positive prime numbers, and we put

\[
\begin{aligned}
(k,p)+\gbr{\frac kp}+\gbr{\frac{2k}{p}}+\gbr{\frac{3k}{p}}\dotsfour+\gbr{\frac{\frac12(p-1)k}{p}}&=L\\
(p,k)+\gbr{\frac pk}+\gbr{\frac{2p}{k}}+\gbr{\frac{3p}{k}}\dotsfour+\gbr{\frac{\frac12(k-1)p}{k}}&=M
\end{aligned}
\]

it is clear from VIII of Article 4 that $L$ and $M$ are always even numbers. But by the theorem of Article 6 we have

\[
\begin{aligned}
L+M=(k,p)+(p,k)+\tfrac14(k-1)(p-1)
\end{aligned}
\]

Thus whenever $\frac14(k-1)(p-1)$ is even, which happens if either both $k,p$ or at least one of them is of the form $4n+1$, $(k,p)$ and $(p,k)$ must necessarily be either both even or both odd. But whenever $\frac14(k-1)(p-1)$ is odd, which happens if both $k,p$ are of the form $4n+3$, one of the numbers $(k,p)$, $(p,k)$ must necessarily be even and the other odd. Thus in the former case the relation of $k$ to $p$ and the relation of $p$ to $k$ (insofar as one is a residue or nonresidue of the other) are identical, and in the latter case opposite.

Q. E. D.

\BookRule{35mm}
